\chapter{热力学与早期宇宙}\label{ux70edux529bux5b66ux4e0eux65e9ux671fux5b87ux5b99}

在二十世纪的漫长岁月里，天文学家和天体物理学家积累了大量证据，表明宇宙在大约137.5亿年前突然诞生——这一事件后来被人们称为"大爆炸"。¹² 天文宇宙起源与演化研究的深入探索，促使物理学与天体物理学逐渐走向融合。热力学与统计力学在我们理解宇宙在大爆炸后不久经历的一系列转变过程中发挥着至关重要的作用。这些转变留下了诸多重要里程碑，天体物理学家正是借助这些里程碑，回溯到宇宙最初始的时刻。早期宇宙尤其适合作为典型案例，用来阐释第六章、第七章和第八章中所阐述的理想经典气体、玻色气体以及费米气体的性质，以及第六章第六节所介绍的化学平衡理论。
\section{大爆炸的观测证据}\label{ux5927ux7206ux70b8ux7684ux89c2ux6d4bux8bc1ux636e}
自埃德温·哈勃在20世纪20年代末发现宇宙正在膨胀以来，大爆炸的观测证据不断增长。自此之后，关于宇宙起源的完整标准模型应运而生。以下三类证据构成了关键的科学依据。
宇宙中几乎每一个星系都在远离其他星系，且它们的退行速度与星系之间的距离之间呈现出近乎线性的依赖关系；详见图~9.1。哈勃率先通过测量邻近星系的距离及其相对于我们银河系的运动速度，发现了这一现象。前者基于标准烛光——在哈勃的案例中，是已知绝对平均亮度的造父变星。后者则基于对光谱线多普勒红移的测量。在利用哈勃太空望远镜进行的最遥远观测中，Ia型超新星被用作标准烛光。这些数据是
¹如需了解大爆炸研究的精彩综述与历史沿革，请参阅韦因伯格所著的《最初三分钟：现代宇宙起源观》（1993年）以及辛格所著的《大爆炸》（2005年）。韦因伯格的《宇宙学》（2008年）提供了极为详尽的技术性概览。本章的编排结构以韦因伯格的著作为基础。美国国家科学院出版社出版的2010年天体物理学十年度调查报告《新世界，新天文学与天体物理学的视野》（New Worlds, New Horizons in Astronomy and Astrophysics），对当前天体物理学领域的最新进展进行了全面概述；详情请访问www.nap.edu。
²稳态宇宙学的倡导者弗雷德·霍伊尔于1950年在一次BBC广播节目中，借机创造了"大爆炸"这一术语。{[}b 令他深感懊恼的是，这个名称很快便广为流传。
被封装在哈勃---弗里德曼关系之中（Friedmann，1922、1924），
\begin{equation}\tag{9.1.1}\label{eq:9.1.1}
\nu = { \frac { d a } { d t } } = H a = { \sqrt { \frac { 8 \pi G u } { 3 c ^{2} } } } a .
\end{equation}
在这里，\(a\) 表示空间中任意两点之间的距离，随着宇宙的膨胀而随时间增长；\(v\) 是退行速度；\(G\) 是万有引力常数；\(c\) 是光速；\(u\) 则是宇宙的能量密度。³ 哈勃参数 \(H\) 是表征宇宙膨胀速率的特征量，其数值大约等于宇宙年龄的倒数。\autoref{eq:9.1.1} 的特定形式假定，正如目前所见，宇宙的能量密度 \(u\) 等于临界值，从而使得
³ 宇宙学和广义相对论的计算通常以等效质量密度 \(\rho = u/c^{2} = T^{00}\) 来表示，其中 \(T\) 是能量-动量张量。天体物理学家通常用多普勒红移因子 \(z = (\lambda/\lambda_{\mathrm{lab}} - 1)\) 来描述长度尺度参数 \(a\)，其中 \(\lambda_{\mathrm{lab}}\) 是某条谱线在实验室中的波长，\(\lambda\) 则是经过红移后的波长。由此可得 \(z = (T/T_{0} - 1)\)，其中 \(T_{0}\) 是当前宇宙微波背景辐射的温度，\(T\) 则是那个时代的光子温度。例如，来自末次散射时期的多普勒红移为
\begin{equation}\tag{9.1.2}\label{eq:9.1.2}
z = \frac { 30 00 \, \mathrm { K } } { 2 . 72 5 \, \mathrm { K } } - 1 \simeq 11 00 ,
\end{equation}
因此，自那时起，宇宙的膨胀幅度已达到 1100 倍。
时空是平坦的。这意味着，宇宙正处在一条刀锋之上——既可能永远持续膨胀，也可能因引力作用而发生坍缩。如需更详尽的技术性综述，请参阅 Börner（2003）和 Weinberg（2008）。哈勃参数的测量值为
\begin{equation}\tag{9.1.3}\label{eq:9.1.3}
H _ { 0 } = 74 . 2 \pm 3 . 6 \, \mathrm { k m \, s ^{- 1} \, M p c ^{- 1} } ,
\end{equation}
其中，Mpc 为兆秒差距，约等于 \(3.26 \times 10^{6}\) 光年；详情请参阅 Freedman 等人（2001）以及 Riess 等人（2009）。
彭齐亚斯和威尔逊（1965）观测到，来自深空的微波辐射噪声几乎均匀且呈各向同性分布，其黑体温度约为3K。这一宇宙微波背景辐射（CMB）很快被确认为大爆炸之后时代遗留下来的黑体辐射。随后，通过气球实验以及美国国家航空航天局"宇宙背景探测器"（COBE）任务开展的太空测量表明，CMB极其均匀且各向同性，平均温度为\(T_{\mathrm{CMB}}=2.725 \pm 0.002 \mathrm{K}\)；参见马瑟等人（1994, 1999）、赖特等人（1994）、菲克森等人（1996），以及图~9.2。美国国家航空航天局"威尔金森微波各向异性探测器"（WMAP）任务对CMB温度的角分布进行了测绘。图~9.3展示了以银河坐标为轴，绘制出的\(\pm 200 \mu \mathrm{K}\) CMB温度变化图。
CMB代表的是在与高能等离子体处于热平衡状态的光子——这种等离子体自宇宙诞生之初一直存在，直到大约在大爆炸后38万年，温度冷却至约3000K。
并由大约72.8\%的暗能量、22.7\%的暗物质以及4.56\%的重子物质（质子和中子）构成。\(^{4}\) 这一比例使得重子数密度 \(n_{B}\) 为 \(0.26\,\mathrm{m}^{-3}\)。在黑体腔内，光子数密度随温度的变化可由\autoref{eq:7.2.33} 给出：
\begin{equation}\tag{9.1.4}\label{eq:9.1.4}
n _ { \gamma } ( T ) = \frac { 2 \zeta ( 3 ) } { \pi ^{2} } \left( \frac { k T } { \hbar c } \right) ^{3} .
\end{equation}
在当前2.725 K的温度下，这将使宇宙微波背景辐射的光子数密度达到 \(n_{\gamma} = 4.10 \times 10^{8}\,\mathrm{m}^{-3}\)，因此当前的重子与光子比为
\begin{equation}\tag{9.1.5}\label{eq:9.1.5}
\eta = \frac { n _ { B } } { n _ { \gamma } } \approx 6 \times 10 ^{- 10} .
\end{equation}
自这两个量随时间 \(a(t)\) 按 \(a^{-3}(t)\) 比例变化以来，\(\eta\) 的比值一直保持不变。正如我们稍后将看到的那样，\(\eta\) 的数值在早期宇宙的热演化过程中发挥了极其重要的作用。\(^{5}\)
在宇宙诞生后的最初几分钟内形成的轻元素 \(^{1}\)H、\(^{2}\)H、\(^{3}\)He、\(^{4}\)He、\(^{7}\)Li 等的相对丰度，都与重子与光子比 \(\eta\) 存在密切关系；参见图~9.4。这一关联最早由乔治·伽莫夫、拉尔夫·阿尔弗和罗伯特·赫尔曼在20世纪40年代末至50年代初加以探索；参见阿尔弗、贝特和伽莫夫（1948年）以及阿尔弗和赫尔曼（1948年、1950年）。\(^{6}\)
\(^{4}\)宇宙的能量含量是通过各类组成成分中所含临界密度的分数来加以参数化的。当前的数值为：暗能量的密度为 \(\Omega_{\Lambda}=0.728\pm0.016\)；重子物质的密度为 \(\Omega_{b}=0.0456\pm0.0016\)；冷暗物质的密度为 \(\Omega_{c}=0.227\pm0.014\)。当前宇宙的年龄，即"回溯时间"，为 \(t_{0}=13.75\pm0.11\times10^{9}\) 年。黑体辐射所占的相对贡献约为 \(6\times10^{-5}\)。这些参数的协和值基于WMAP 7年数据，并已列于Komatsu等人（2010）的表格中。暗能量是导致宇宙加速膨胀的原因。在宇宙早期，能量密度的分布差异巨大，因为它们随膨胀参数 \(a\) 的变化而呈现出不同的比例关系。在再结合时期，各成分的占比分别为：暗物质占63\%，重子物质占12\%，相对论性辐射（光子和中微子）占25\%。在最初的数百万年里，相对论性粒子对宇宙能量的贡献最为显著。根据表9.2中给出的光子：中微子：电子的比例——2:21/4:7/2——计算得出，光子占能量的18.6\%，中微子和反中微子占48.8\%，电子与正电子占32.5\%。尽管目前暗能量已成为宇宙能量密度的主要组成部分，但在宇宙早期的演化过程中，它所起的作用却十分有限。冷暗物质在"黑暗时代"结束时——也就是在最后一次散射后1亿至2亿年——对于第一代恒星和星系的形成至关重要。
\(^{5}\)此处采用的正确计量方式是重子数密度与光子熵密度之比；然而，由于CMB的光子熵密度和数密度均按 \(T^{3}\) 按照比例变化，因此该比值通常以数密度之比的形式进行表述。
\$\^{}\{6\}\$1948年，乔治·伽莫夫和拉尔夫·阿尔弗提出了一种模型，用于解释在炽热且不断膨胀的原始汤中，质子、中子和电子发生核合成的过程。阿尔弗和伽莫夫将这种物质称为"ylem"。为了解释宇宙中现今存在的 \(^{4}\)He丰度，阿尔弗和赫尔曼（1950）提出了一个大约为 \(10^{-9}\) 的重子与光子比，并预测了当前宇宙微波背景辐射的温度约为5K。伽莫夫将他朋友汉斯·贝特的名字作为第二作者，加入到阿尔弗、贝特和伽莫夫（1948）的论文中，以此谐音取自希腊字母表。这篇论文发表于4月1日——或许并非巧合。参见阿尔弗和赫尔曼（2001）、温伯格（1993）以及辛格（2005）。
这些标记，为我们提供了关于宇宙最早时刻属性与行为的有力证据。
从宇宙诞生的最初时刻，直至38万年后进入再结合时期，宇宙等离子体通过汤姆逊散射与黑体辐射处于热平衡状态。由于带电粒子的密度极高，光子散射的平均自由程远短于宇宙在膨胀与冷却过程中温度变化的时间尺度，因此等离子体始终与光子保持热平衡。在宇宙膨胀的前几十万年里，宇宙的能量密度主要由光子及其他相对论性粒子所主导。这是因为黑体辐射的能量密度与 \(a^{-4}\) 成正比，而非相对论性物质的能量密度则与 \(a^{-3}\) 成正比。图~9.5 和表9.1展示了黑体光子的温度随宇宙年龄的变化情况。
在最初的百分之一秒内，宇宙从其单一的起始状态迅速膨胀并冷却，温度降至约 \(10^{11}\) K。自此以后，宇宙的物理规律主要由弱相互作用和电磁相互作用所支配。由于重子与光子的比例极小，且温度已降至
\section{相对论电子、正电子与中微子}\label{ux76f8ux5bf9ux8bbaux7535ux5b50ux6b63ux7535ux5b50ux4e0eux4e2dux5faeux5b50}
在宇宙最早期的时刻，温度高到足以产生多种相对论性粒子和反粒子。如果 \(kT \gg mc^2\)，那么每个具有质量 \(m\) 的粒子-反粒子对都可以通过光子-光子相互作用而产生。在这样的温度下，几乎所有的生成粒子都将满足由相对论极限所描述的能量-动量关系，即
\begin{equation}\tag{9.3.1}\label{eq:9.3.1}
\varepsilon _ { \mathbf { k } } \approx \hbar c k ,
\end{equation}
\(^{7}\)在时间 \(t=0.01\,\mathrm{s}\) 之前，由于强相互作用的粒子和反粒子的产生，分析工作变得更加复杂。更早的时期，当温度高于 \(kT\approx300\,\mathrm{MeV}\)（\(T=4\times10^{12}\,\mathrm{K}\)）时，强子已分解为一种具有强相互作用的相对论性夸克-胶子等离子体。布鲁克海文国家实验室的相对论重离子对撞机成功创造了迄今为止在实验室中所制造出的温度最高的物质——等离子体，其温度高达 \(T=4\times10^{12}\,\mathrm{K}\)；参见 Adare 等人（2010）。
\(^{8}\)关于重子生成的精确机制（即非零重子与光子比 \(\eta\)）仍存在争议。正如萨哈罗夫（1967）所指出的那样，这一过程需要具备三方面条件：重子数不守恒、C 反对称性与 CP 反对称性破缺，以及偏离平衡态。这些条件在宇宙最早的时刻均已得到满足（远早于我们在此处所研究的时间尺度），但至今仍未形成一套共识性的理论，能够解释重子不对称性，并使其大小接近于观测值 \(\eta=6\times10^{-10}\)。
其中 \(\hbar k\) 表示动量的大小。该关系适用于光子、中微子、反中微子、电子以及正电子。电子-正电子对形成的阈值为 \(m_{e}c^{2}/k=5.9\times10^{9}\) K。中微子质量极小，因此我们可以安全地假设它们处于相对论性状态。相对论性色散\autoref{eq:9.3.1} 为所有相对论性粒子种类提供了几乎相同的态密度：
\begin{equation}\tag{9.3.2}\label{eq:9.3.2}
a ( \varepsilon ) = \frac { g _ { s } } { ( 2 \pi ) ^{3} } \int \delta ( \varepsilon - \varepsilon _ { k } ) d \mathbf { k } = \frac { 4 \pi g _ { s } } { ( 2 \pi ) ^{3} } \int _ { 0 } ^{\infty} k ^{2} \delta ( \varepsilon - \hbar c k ) d \mathbf { k } = \frac { g _ { s } \varepsilon ^{2} } { 2 \pi ^{2} ( \hbar c ) ^{3} } ,
\end{equation}
其中 \(\mathbf{g}_{s}\) 表示自旋简并度。光子的自旋简并度为 \(\mathbf{g}_{s}=2\)（左旋和右旋圆偏振）。其他粒子均为自旋-\(\frac{1}{2}\) 费米子。电子和正电子的自旋简并度为 \(\mathbf{g}_{s}=2\)，而中微子和反中微子的自旋简并度则为 \(\mathbf{g}_{s}=1\)，因为所有中微子都具有左手螺旋自旋。
在这一时期，由于宇宙的电荷中性以及重子与光子比 \(\eta\) 的数值较小，电子和正电子的数目密度几乎相等，因此它们的化学势都相当低。假设宇宙的轻子总数也相对较低，那么中微子和反中微子同样适用这一规律。如第 7.3 节所阐述的那样，光子的化学势恰好为零。对于由费米子（+）或玻色子（-）组成的相对论性气体，其化学势、压力、数目密度、能量密度以及熵密度均可表示为：
\begin{equation}\tag{9.3.3}\label{eq:9.3.3}
P ( T ) = \pm k T \int a ( \varepsilon ) \ln ( 1 \pm e ^{- \beta \varepsilon} ) d \varepsilon = \frac { g _ { s } \left( k T \right) ^{4} } { 2 \pi ^{2} \left( \hbar c \right) ^{3} } \int _ { 0 } ^{\infty} x ^{2} \ln \left( 1 \pm e ^{- x} \right) d x ,
\end{equation}
\begin{equation}\tag{9.3.4}\label{eq:9.3.4}
n ( T ) = \int a ( \varepsilon ) \frac { 1 } { e ^{\beta \varepsilon} \pm 1 } d \varepsilon = \frac { g s } { 2 \pi ^{2} } \left( \frac { k T } { \hbar c } \right) ^{3} \int _ { 0 } ^{\infty} \frac { x ^{2} } { e ^{x} \pm 1 } d x ,
\end{equation}
\begin{equation}\tag{9.3.5}\label{eq:9.3.5}
u ( T ) = \int a ( \varepsilon ) \frac { \varepsilon } { e ^{\beta \varepsilon} \pm 1 } d \varepsilon = \frac { g _ { s } \left( k T \right) ^{4} } { 2 \pi ^{2} \left( \hbar c \right) ^{3} } \int _ { 0 } ^{\infty} \frac { x ^{3} } { e ^{x} \pm 1 } d x ,
\end{equation}
\begin{equation}\tag{9.3.6}\label{eq:9.3.6}
s ( T ) = \left( \frac { \partial P } { \partial T } \right) _ { \mu } = \frac { 2 g _ { s } k } { \pi ^{2} } \left( \frac { k T } { \hbar c } \right) ^{3} \int _ { 0 } ^{\infty} x ^{2} \ln \left( 1 \pm e ^{- x} \right) d x .
\end{equation}
\(^{9}\)已知所有三种中微子族均因中微子振荡观测而呈现出微小（但非零）的质量。其中，电子中微子的质量可能最轻，实验上限为 \(m_{\nu_{e}}c^{2}<2.2\) eV。WMAP 拍摄的宇宙微波背景辐射角涨落分布，为中微子总质量之和设定了上限，即 \(\sum m_{\nu}c^{2}<0.58\) eV。因此，我们可以安全地假设：在早期宇宙中，所有中微子种类的质量都远低于 \(kT\) 的数值。
利用附录D中给出的玻色积分值，我们得到了黑体光子的以下表达式：
\begin{equation}\tag{9.3.7}\label{eq:9.3.7}
P _ { \gamma } ( T ) = \frac { \pi ^{2} } { 45 } \frac { ( k T ) ^{4} } { ( \hbar c ) ^{3} } ,
\end{equation}
\begin{equation}\tag{9.3.8}\label{eq:9.3.8}
n _ { \gamma } ( T ) = \frac { 2 \zeta ( 3 ) } { \pi ^{2} } \left( \frac { k T } { \hbar c } \right) ^{3} ,
\end{equation}
\begin{equation}\tag{9.3.9}\label{eq:9.3.9}
u _ { \gamma } ( T ) = \frac { \pi ^{2} } { 15 } \frac { ( k T ) ^{4} } { ( \hbar c ) ^{3} } ,
\end{equation}
\begin{equation}\tag{9.3.10}\label{eq:9.3.10}
s _ { \gamma } ( T ) = \frac { 4 \pi ^{2} k } { 45 } \left( \frac { k T } { \hbar c } \right) ^{3} .
\end{equation}
所有μ=0 的相对论性粒子，在压力、能量密度等物理量的温度依赖关系上，与光子完全相同；而费米积分与玻色积分的区别仅在于一个常数因子：
\begin{equation}\tag{9.3.11}\label{eq:9.3.11}
\int _ { 0 } ^{\infty} { \frac { x ^{n - 1} } { e ^{x} + 1 } } d x = \left( 1 - { \frac { 1 } { 2 ^{n - 1} } } \right) \int _ { 0 } ^{\infty} { \frac { x ^{n - 1} } { e ^{x} - 1 } } d x ;
\end{equation}
参见附录D和附录E。压力、能量密度和熵密度的贡献，源自对自旋简并度、粒子数与反粒子数的统计，以及对不同费米/玻色因子（费米子为1，费米子为7/8）的考量。光子、三代中微子（电子中微子、缪子中微子、τ中微子及其反粒子），以及电子和正电子，以表9.2所示的比例，共同构成了总压力、粒子数密度、能量密度和熵密度。此处的计数方式与文献中通常采用的方式一致——以光子每种自旋态的贡献为基准。粒子数密度的贡献保持不变，只是费米/玻色因子现为3/4。
表9.2 相对论性粒子对压力、能量密度和熵密度的贡献
\begin{longtable}[]{@{}
  >{\raggedright\arraybackslash}p{(\columnwidth - 8\tabcolsep) * \real{0.2000}}
  >{\raggedright\arraybackslash}p{(\columnwidth - 8\tabcolsep) * \real{0.2000}}
  >{\raggedright\arraybackslash}p{(\columnwidth - 8\tabcolsep) * \real{0.2000}}
  >{\raggedright\arraybackslash}p{(\columnwidth - 8\tabcolsep) * \real{0.2000}}
  >{\raggedright\arraybackslash}p{(\columnwidth - 8\tabcolsep) * \real{0.2000}}@{}}
\toprule\noalign{}
\begin{minipage}[b]{\linewidth}\raggedright
粒子
\end{minipage} & \begin{minipage}[b]{\linewidth}\raggedright
费米/玻色因子
\end{minipage} & \begin{minipage}[b]{\linewidth}\raggedright
自旋简并度
\end{minipage} & \begin{minipage}[b]{\linewidth}\raggedright
粒子种类数
\end{minipage} & \begin{minipage}[b]{\linewidth}\raggedright
\(2P_{\rm total}/P_{\gamma}\)
\end{minipage} \\
\midrule\noalign{}
\endhead
\bottomrule\noalign{}
\endlastfoot
\(\gamma\) & \(1\) & \(2\) & \(1\) & \(2\) \\
\(\nu_{e}, \nu_{\mu}, \nu_{\tau}\) & \(\frac{7}{8}\) & \(1\) & \(3\) & \(\frac{21}{8}\) \\
\(\bar{\nu}_{e}, \bar{\nu}_{\mu}, \bar{\nu}_{\tau}\) & \(\frac{7}{8}\) & \(1\) & \(3\) & \(\frac{21}{8}\) \\
\(e^{-}\) & \(\frac{7}{8}\) & \(2\) & \(1\) & \(\frac{7}{4}\) \\
\(e^{+}\) & \(\frac{7}{8}\) & \(2\) & \(1\) & \(\frac{7}{4}\) \\
\end{longtable}
最终的总和如下：
\begin{equation}\tag{9.3.12}\label{eq:9.3.12}
P _ { \mathrm { t o t a l } } ( T ) = \left( 2 + \frac { 21 } { 4 } + \frac { 7 } { 2 } \right) \frac { P _ { \gamma } ( T ) } { 2 } = \left( \frac { 43 } { 8 } \right) P _ { \gamma } ( T ) ,
\end{equation}
\begin{equation}\tag{9.3.13}\label{eq:9.3.13}
u _ { \mathrm { t o t a l } } ( T ) = \left( 2 + \frac { 21 } { 4 } + \frac { 7 } { 2 } \right) \frac { u _ { \gamma } ( T ) } { 2 } = \left( \frac { 43 } { 8 } \right) u _ { \gamma } ( T ) ,
\end{equation}
\begin{equation}\tag{9.3.14}\label{eq:9.3.14}
s _ { \mathrm { t o t a l } } ( T ) = \left( 2 + \frac { 21 } { 4 } + \frac { 7 } { 2 } \right) \frac { s _ { \gamma } ( T ) } { 2 } = \left( \frac { 43 } { 8 } \right) s _ { \gamma } ( T ) ,
\end{equation}
\begin{equation}\tag{9.3.15}\label{eq:9.3.15}
n _ { \mathrm { t o t a l } } ( T ) = \left( 2 + \frac { 9 } { 2 } + 3 \right) \frac { n _ { \gamma } ( T ) } { 2 } = \left( \frac { 19 } { 4 } \right) n _ { \gamma } ( T ) .
\end{equation}
在这一时期，宇宙的密度足够高，因此弱相互作用和电磁相互作用的速率使所有这些粒子都与彼此处于热平衡状态。因此，随着宇宙以绝热方式膨胀，当体积从初始值 \(a_{0}^{3}\) 扩展至最终体积 \(a_{1}^{3}\) 时，具有线性尺寸 \(a\) 的共动体积中的熵保持不变：
\begin{equation}\tag{9.3.16}\label{eq:9.3.16}
s _ { \mathrm { t o t a l } } ( T _ { 0 } ) a _ { 0 } ^{3} = s _ { \mathrm { t o t a l } } ( T _ { 1 } ) a _ { 1 } ^{3} .
\end{equation}
由于熵密度与 \(T^{3}\) 成正比，因此在时间 \(t\) 时的温度与长度尺度之间存在如下关系：
\begin{equation}\tag{9.3.17}\label{eq:9.3.17}
T ( t ) a ( t ) = \mathrm { c o n s t . }
\end{equation}
这与自由膨胀的光子气体所适用的关系相同，详见\autoref{eq:9.3.17}，只不过在这里，该关系源自一个绝热平衡过程。根据\autoref{eq:9.1.1} 和 \autoref{eq:9.3.13}，在这一时期，宇宙的温度随宇宙年龄 \(t\) 的变化而变化，其表达式为：
\begin{equation}\tag{9.3.18}\label{eq:9.3.18}
T ( t ) = 10 ^{10} \, \mathrm { K } \sqrt { \frac { 0 . 99 2 \mathrm { s } } { t } } \, ;
\end{equation}
参见问题 9.1。
\section{中子分数}\label{ux4e2dux5b50ux5206ux6570}
在宇宙诞生后的第一秒内，当 \(T > 10^{10}\) K、且质子与中子尚未结合成原子核时，弱相互作用使游离的中子与质子在热"β-平衡"状态下彼此相持，并通过以下过程与光子、中微子、电子以及正电子保持平衡：
\begin{equation}\tag{9.4.1}\label{eq:9.4.1}
n + \nu _ { e } \rightleftharpoons p + e ^{-} + \gamma ,
\end{equation}
\begin{equation}\tag{9.4.2}\label{eq:9.4.2}
n + e ^{+} \rightleftharpoons p + \bar { \nu } + \gamma ,
\end{equation}
\begin{equation}\tag{9.4.3}\label{eq:9.4.3}
n \rightleftharpoons p + e ^{-} + \bar { \nu } + \gamma .
\end{equation}
我们可以将此视为一种化学平衡过程，正如第 6.6 节所描述的那样。由于光子、电子、正电子、中微子以及反中微子的化学势均为零，因此在平衡状态下，中子与质子的化学势必须相等：
\begin{equation}\tag{9.4.4}\label{eq:9.4.4}
\mu _ { n } = \mu _ { p } .
\end{equation}
在这些温度 (\(\approx 10^{11}\) K) 和密度 (\(\approx 10^{32}\) \(m^{-3}\)) 下，质子与中子可以被视作经典非相对论的理想气体。依据\autoref{eq:6.6.5}，自旋 \(\frac{1}{2}\) 的质子与中子的化学势分别为：
\begin{equation}\tag{9.4.5}\label{eq:9.4.5}
\mu _ { p } = m _ { p } c ^{2} + k T \mathrm { l n } \left( n _ { p } \lambda _ { p } ^{3} \right) - k T \mathrm { l n } 2 ,
\end{equation}
\begin{equation}\tag{9.4.6}\label{eq:9.4.6}
\mu _ { n } = m _ { n } c ^{2} + k T \ln \left( n _ { n } \lambda _ { n } ^{3} \right) - k T \ln 2 .
\end{equation}
其中 \(\lambda(=h/\sqrt{2\pi mkT})\) 是热德布罗意波长。中子的静止能量比质子的静止能量高出
\begin{equation}\tag{9.4.7}\label{eq:9.4.7}
m _ { n } c ^{2} - m _ { p } c ^{2} = \Delta \varepsilon = 1 . 29 3 \mathrm { M e V } .
\end{equation}
忽略\autoref{eq:9.3.3} 和\autoref{eq:9.3.3} 中热德布罗意波长的微小质量差异，可得
\begin{equation}\tag{9.4.8}\label{eq:9.4.8}
n _ { n } = n _ { p } e ^{- \beta \Delta \varepsilon} .
\end{equation}
重子数密度等于中子数密度与质子数密度之和
\begin{equation}\tag{9.4.9}\label{eq:9.4.9}
n _ { B } = n _ { n } + n _ { p } ,
\end{equation}
因此，平衡状态下的中子分数可由下式给出：
\begin{equation}\tag{9.4.10}\label{eq:9.4.10}
q = \frac { n _ { n } } { n _ { B } } = \frac { 1 } { e ^{\beta \Delta \varepsilon} + 1 } .
\end{equation}
质量差导致的跨临界温度为 \(T_{np} = \Delta\varepsilon/k \approx 1.50 \times 10^{10}\,\mathrm{K}\)，因此，在 \(T=10^{11}\,\mathrm{K}\) 时，中子比例从 46\% 降至 \(T=9\times10^{9}\,\mathrm{K}\) 时的 16\%，而 \(t_{1} \approx 1\,\mathrm{秒}\)。随着温度下降至低于 \(10^{10}\,\mathrm{K}\)（\(kT=0.86\,\mathrm{MeV}\)），弱相互作用率开始远低于宇宙的冷却速率，因此重子迅速脱离与中微子的平衡状态。自此之后，中子开始发生β衰变，其自然放射性衰变寿命为 \(\tau_{n}=886\,\mathrm{秒}\)，因此中子比例
呈指数级衰减：
\begin{equation}\tag{9.4.11}\label{eq:9.4.11}
q \approx 0 . 16 \mathrm { e x p } \left( \frac { - ( t - t _ { 1 } ) } { \tau _ { n } } \right)  \mathrm { f o r } \; t > t _ { 1 } = 1 \mathrm { s } .
\end{equation}
在核合成发生之时，大约 3.7 分钟后，中子比例已降至 \(q \approx 0.12\)。此时，剩余的中子与质子结合形成氘核及其他轻核。有关核合成的详细讨论，请参见第 9.7 节。
\section{正电子与电子的湮灭}\label{ux6b63ux7535ux5b50ux4e0eux7535ux5b50ux7684ux6e6eux706d}
大爆炸后约一秒钟，温度接近产生电子-正电子对的跨临界温度 \(T_{e}\)：
\begin{equation}\tag{9.5.1}\label{eq:9.5.1}
k T _ { e } = m _ { e } c ^{2} = 0 . 51 1 \, \mathrm { M e V } ,
\end{equation}
当 \(T_{e}=5.93 \times 10^{9}\,\mathrm{K}\) 时。随着温度下降至低于 \(T_{e}\)，\(e^{+}e^{-}\) 对的生成速率开始低于对湮灭的速率。电子的完整相对论性色散关系为
\begin{equation}\tag{9.5.2}\label{eq:9.5.2}
\varepsilon _ { k } = \sqrt { ( \hbar c k ) ^{2} + ( m _ { e } c ^{2} ) ^{2} } \; ,
\end{equation}
该公式给出了态密度
\begin{equation}\tag{9.5.3}\label{eq:9.5.3}
a _ { e } ( \varepsilon ) = \frac { 8 \pi } { ( 2 \pi ) ^{3} } \int _ { 0 } ^{\infty} k ^{2} \delta ( \varepsilon - \varepsilon _ { k } ) d k = \frac { \varepsilon \sqrt { \varepsilon ^{2} - ( m _ { e } c ^{2} ) ^{2} } } { \pi ^{2} ( \hbar c ) ^{3} }  \mathrm { f o r } \; \varepsilon \ge m _ { e } c ^{2} \, .
\end{equation}
由于电子和正电子通过反应与黑体光子处于平衡状态，
\begin{equation}\tag{9.5.4}\label{eq:9.5.4}
e ^{+} + e ^{-} \rightleftharpoons \gamma + \gamma ,
\end{equation}
根据\autoref{eq:6.6.3}，平衡方程表明各物种的化学势之间存在如下关系：
\begin{equation}\tag{9.5.5}\label{eq:9.5.5}
\mu _ { - } + \mu _ { + } = 2 \mu _ { \gamma } = 0 .
\end{equation}
电子数密度与光子数密度之比随后为
\begin{equation}\tag{9.5.6}\label{eq:9.5.6}
\frac { n _ { - } } { n _ { \gamma } } = \frac { 1 } { 2 \zeta ( 3 ) } \int _ { \beta m _ { e } c ^{2} } ^{\infty} \frac { x \sqrt { x ^{2} - ( \beta m _ { e } c ^{2} ) ^{2} } } { e ^{x} e ^{- \beta \mu -} + 1 } d x ,
\end{equation}
而正电子数密度比则为
\begin{equation}\tag{9.5.7}\label{eq:9.5.7}
\frac { n _ { + } } { n _ { \gamma } } = \frac { 1 } { 2 \zeta ( 3 ) } \int _ { \beta m _ { e } c ^{2} } ^{\infty} \frac { x \sqrt { x ^{2} - ( \beta m _ { e } c ^{2} ) ^{2} } } { e ^{x} e ^{\beta \mu -} + 1 } d x ;
\end{equation}
参见\autoref{eq:9.1.5}。随着宇宙降温，电子与正电子与黑体光子达到平衡的状态被打破。
最终，所有正电子都被湮灭，仅留下目前仍然存在的电子。宇宙的电荷中性要求电子数密度与正电子数密度之差等于质子数密度，即 \((1-q)n_{B}\)，其中 \(q\) 的定义见第 9.4 节；因此
\begin{equation}\tag{9.5.8}\label{eq:9.5.8}
\frac { ( n _ { - } - n _ { + } ) } { n _ { \gamma } } = \frac { \sinh ( \beta \mu _ { - } ) } { 2 \zeta ( 3 ) } \int _ { \beta m _ { e } c ^{2} } ^{\infty} \frac { x \sqrt { x ^{2} - ( \beta m _ { e } c ^{2} ) ^{2} } } { \cosh ( x ) + \cosh ( \beta \mu _ { - } ) } d x = ( 1 - q ) \eta .
\end{equation}
我们可以利用\autoref{eq:9.5.8} 以数值方法确定电子化学势随温度的变化，然后将该值代入\autoref{eq:9.5.6} 和 \autoref{eq:9.5.7} 中，以计算电子和正电子的数密度；请参见图~9.6。
起初，随着温度降至电子-正电子对阈值以下，电子和正电子的密度均按\(\exp(-\beta m_{e}c^{2})\)的比例下降，但它们
而在温度\(T_{1}\)时，由于此时几乎所有的电子和正电子均已湮灭，因此该温度下的熵仅由光子提供：
\begin{equation}\tag{9.5.9}\label{eq:9.5.9}
S ( T _ { 1 } ) = s _ { \gamma } ( T _ { 1 } ) a _ { 1 } ^{3} .
\end{equation}
在绝热膨胀过程中，熵的守恒关系将初始温度与最终温度联系起来，具体为
\begin{equation}\tag{9.5.10}\label{eq:9.5.10}
\left( \frac { 11 } { 4 } \right) ^{1 / 3} T _ { 0 } a _ { 0 } = T _ { 1 } a _ { 1 } .
\end{equation}
本质上，湮灭电子与正电子的熵被传递给了光子。由于中微子与光子在电子-正电子湮灭之前温度相等，且在湮灭过程中中微子得以自由膨胀，因此\ldots\ldots{}
在湮灭时期，中微子的温度下降得比光子的温度更为显著：
\begin{equation}\tag{9.5.11}\label{eq:9.5.11}
T _ { \nu 1 } = ( 4 / 11 ) ^{1 / 3} \, T _ { 1 } .
\end{equation}
在\(e^{+}e^{-}\)湮灭之后，中微子和光子的温度均按照 \autoref{eq:9.3.17} 演化，因此，宇宙大爆炸遗留中微子的当前温度应为
\begin{equation}\tag{9.5.12}\label{eq:9.5.12}
T _ { \nu } = ( 4 / 11 ) ^{1 / 3} \, T _ { \mathrm { C M B } } \simeq 1 . 94 5 \, \mathrm { K } .
\end{equation}
对宇宙中微子背景的测量，将为标准模型的宇宙大爆炸理论提供一项极佳的额外验证手段，但目前我们尚无可行的方法来探测这些能量极其低的中微子。\(^{10}\)
\section{原始核合成}\label{ux539fux59cbux6838ux5408ux6210}
除氢以外的轻核是在宇宙大爆炸后3到4分钟之间形成的，当时温度已降至约\(10^{9}\) K。在此之前，高温黑体辐射会迅速使偶然形成的氘核发生光解离。从质子和中子中形成轻核的第一步是生成氘，因为在这些密度下，核的形成速率主要由双体碰撞所主导。一旦氘核形成，大部分此类核会在随后的一系列双体碰撞中，与剩余的质子、中子以及彼此发生反应，迅速转化为氦以及其他更稳定的轻核。正如第9.4节所讨论的那样，此时质子与中子的混合比例约为\(q=12%
\)中子，\(1-q=88%
\)质子。到了\(t\approx3\)分钟时，温度已降至\(T\approx10^{9}\) K，因此质子和中子得以通过这一过程开始相互结合，形成氘核。
\begin{equation}\tag{9.7.1}\label{eq:9.7.1}
p + n \rightleftharpoons d + \gamma \, .
\end{equation}
该反应的化学平衡关系，见\autoref{eq:6.6.3}，为
\begin{equation}\tag{9.7.2}\label{eq:9.7.2}
\mu _ { p } + \mu _ { n } = \mu _ { d } ,
\end{equation}
因为黑体光子的化学势为零。在这些温度和密度下，质子、中子和氘核可以被视为经典理想气体。
质子和中子均为自旋-\(\frac{1}{2}\)粒子，因此它们各自具有两种自旋态；而氘核则自旋为1，拥有三种自旋态：
\begin{equation}\tag{9.7.3}\label{eq:9.7.3}
\mu _ { p } = m _ { p } c ^{2} + k T \mathrm { l n } \left( n _ { p } \lambda _ { p } ^{3} \right) - k T \mathrm { l n } 2 ,
\end{equation}
\begin{equation}\tag{9.7.4}\label{eq:9.7.4}
\mu _ { n } = m _ { n } c ^{2} + k T \ln \left( n _ { n } \lambda _ { n } ^{3} \right) - k T \ln 2 ,
\end{equation}
\begin{equation}\tag{9.7.5}\label{eq:9.7.5}
\mu _ { d } = m _ { d } c ^{2} + k T \ln \left( n _ { d } \lambda _ { d } ^{3} \right) - k T \ln 3 .
\end{equation}
氘核的结合能为 \(\varepsilon_{b}=m_{p}c^{2}+m_{n}c^{2}-m_{d}c^{2}=2.20\,\mathrm{MeV}\)。由于氘核的质量大约是质子或中子的两倍，因此氘核的数密度可由以下公式给出：
\begin{equation}\tag{9.7.6}\label{eq:9.7.6}
n _ { d } = \frac { 3 } { 4 } n _ { p } n _ { n } \frac { \lambda _ { p } ^{3} \lambda _ { n } ^{3} } { \lambda _ { d } ^{3} } e ^{\beta \varepsilon _ { b} } \approx \frac { 3 } { \sqrt { 2 } } n _ { p } n _ { n } \lambda _ { p } ^{3} e ^{\beta \varepsilon _ { b} } .
\end{equation}
重子的总数密度由重子与光子的比值决定：\(\eta\)：\(n_{B}=\eta n_{\gamma}=n_{p}+n_{n}+2n_{d}\)。中子数密度为 \(qn_{B}=n_{n}+n_{d}\)，因此氘核分数可由下式给出：
\begin{equation}\tag{9.7.7}\label{eq:9.7.7}
f _ { d } = \frac { n _ { d } } { n _ { B } } = ( 1 - q - f _ { d } ) ( q - f _ { d } ) s ,
\end{equation}
其中参数 \(s\) 为
\begin{equation}\tag{9.7.8}\label{eq:9.7.8}
s = \frac { 12 \zeta ( 3 ) } { \sqrt { \pi } } \left( \frac { k T } { m _ { p } c ^{2} } \right) ^{3 / 2} \eta e ^{\beta \varepsilon _ { b} } ;
\end{equation}
另请参见\autoref{eq:9.1.5}。\autoref{eq:9.7.7} 与第9.8节将要讨论的氢原子电离的萨哈方程相似，并且其解为
\begin{equation}\tag{9.7.9}\label{eq:9.7.9}
f _ { d } = \frac { 1 + s - \sqrt { ( 1 + s ) ^{2} - 4 s ^{2} q ( 1 - q ) } } { 2 s } .
\end{equation}
在高温条件下，\(s\) 较小，\(f_{d} \approx q(1-q)s\)；而在低温条件下，\(s\) 则较大，\(f_{d} \approx q\)，也就是说，所有中子都被束缚成氘核。图~9.7展示了氘核分数随温度的变化曲线。由于重子与光子的比值 \(\eta\) 和 \(\varepsilon_{b}/m_{p}c^{2}\) 的数值较小，这使得核合成过程被推迟，直到温度降至
\begin{equation}\tag{9.7.10}\label{eq:9.7.10}
k T _ { n } \approx \frac { \varepsilon _ { b } } { \ln \left( \frac { 1 } { \eta } \left( \frac { m _ { p } c ^{2} } { \varepsilon _ { b } } \right) ^{3 / 2} \right) } \, ,
\end{equation}
为中子分数衰减至 \(q=0.12\) 提供了时间。
\begin{equation}\tag{9.7.11}\label{eq:9.7.11}
d + d \rightarrow ^{3} \mathrm { H e } + n + \gamma ,
\end{equation}
\begin{equation}\tag{9.7.12}\label{eq:9.7.12}
d + \mathrm { ^ 3 H } \rightarrow \mathrm { ^ 4 H e } + n + \gamma ,
\end{equation}
\begin{equation}\tag{9.7.13}\label{eq:9.7.13}
d + \mathrm { ^ 3 H e } \rightarrow \mathrm { ^ 4 H e } + p + \gamma ,
\end{equation}
结果是，几乎所有氘核都被转化为极其稳定的同位素\(^{4}\)He，同时还有少量其他轻核。由于每个\(^{4}\)He原子核由两个质子和两个中子组成，因此氦的质量分数为\(2q=24\)\%，质子质量分数为\(1-2q=76\)\%。完整的计算过程涉及非平衡效应，这些效应通过各核相互作用的速率方程进行建模，包括对更重同位素的相互作用；不过，这仅会略微改变\(^{4}\)He的预测浓度。关于\(^{11}\)，请参阅温伯格（2008）的研究。令人惊讶的是，理论上的最大不确定性竟在于\autoref{eq:9.4.8} 中中子的放射性衰变时间的不确定性；详情请参阅科皮、施拉姆和特纳（1997）。这类计算最早由伽莫夫、阿尔弗和赫尔曼在20世纪40年代末至50年代初完成。基于当前的氦及其他轻元素的含量
在宇宙中，阿尔弗和赫尔曼早在彭齐亚斯和威尔逊发现之前十年就预言了5K的宇宙微波背景辐射；请参阅脚注6。
\section{再结合}\label{ux518dux7ed3ux5408}
在最初的几分钟内完成核合成之后，宇宙继续冷却，原子核与电子仍以普通等离子体的形式存在，并与光子处于热平衡状态。直到数十万年之后，温度才降至低于原子电离能——大约几电子伏特，而这一能量足以使原子核俘获电子并形成原子。氢是最后形成的中性粒子，因为它具有最低的里德伯电离能（1Ry = \(m_{e}e^{4}/8\epsilon_{0}^{2}h^{2} = 13.6057\)eV）。乍一看，人们可能会认为，当温度降至里德伯/克数等于158,000 K时，原子就会形成；然而，正如我们接下来将要看到的那样，每一个质子所携带的大量光子，却将再结合的时间推迟到了\(T \approx 3000\)K。一旦所有电子与质子都形成了中性氢原子，宇宙便因自由电子最后一次散射辐射而变得透明。这些CMB黑体光子瞬间获得了自由传播的权利，自那时起便一直未受散射地传播至今。
再结合反应（即氢光电离反应的逆反应）是
\begin{equation}\tag{9.8.1}\label{eq:9.8.1}
p + e \rightleftharpoons \mathrm { H } + \gamma ,
\end{equation}
因此，第6.6节中的化学平衡关系给出了
\begin{equation}\tag{9.8.2}\label{eq:9.8.2}
\mu _ { p } + \mu _ { e } = \mu _ { \mathrm { H } }
\end{equation}
因为，再次强调，黑体光子的化学势为零。在这一时期所处的温度与密度下（数千开尔文，且每立方米仅约\(10^{9}\)个原子），电子、质子和氢原子均可被视为经典理想气体，其结果是
\begin{equation}\tag{9.8.3}\label{eq:9.8.3}
\mu _ { p } = m _ { p } c ^{2} + k T \ln ( n _ { p } \lambda _ { p } ^{3} ) - k T \ln 2 ,
\end{equation}
\begin{equation}\tag{9.8.4}\label{eq:9.8.4}
\mu _ { e } = m _ { e } c ^{2} + k T \ln ( n _ { e } \lambda _ { e } ^{3} ) - k T \ln 2 ,
\end{equation}
\begin{equation}\tag{9.8.5}\label{eq:9.8.5}
\mu _ { \mathrm { H } } = m _ { \mathrm { H } } c ^{2} + k T \ln ( n _ { \mathrm { H } } \lambda _ { \mathrm { H } } ^{3} ) - k T \ln 4 .
\end{equation}
氢的结合能为 \(m_{p}c^{2} + m_{e}c^{2} - m_{\mathrm{H}}c^{2} = 1\) Ry。由此，平衡条件（6.6.3）与理想气体化学势（6.6.5）共同给出了三种粒子数密度之间的一种简单关系：
\begin{equation}\tag{9.8.6}\label{eq:9.8.6}
n _ { \mathrm { H } } = n _ { p } n _ { e } \lambda _ { e } ^{3} e ^{\beta \mathrm { R y} } ,
\end{equation}
其中
\begin{equation}\tag{9.8.7}\label{eq:9.8.7}
\lambda _ { e } = \frac { h } { \sqrt { 2 \pi \, m _ { e } k T } }
\end{equation}
是电子热德布罗意波长。由于电荷中性，自由电子和质子的数密度相同：
\begin{equation}\tag{9.8.8}\label{eq:9.8.8}
n _ { e } = n _ { p } .
\end{equation}
核合成后剩余的质子或以自由形式存在，或结合成氢原子，因此
\begin{equation}\tag{9.8.9}\label{eq:9.8.9}
n _ { p } + n _ { \mathrm { H } } = ( 1 - 2 q ) n _ { B } = ( 1 - 2 q ) \eta n _ { \gamma } .
\end{equation}
将\autoref{eq:9.8.3}、\autoref{eq:9.8.4}、\autoref{eq:9.8.6} 和 \autoref{eq:9.8.8} 综合起来，并利用\autoref{eq:9.1.5}，可得中性氢丰度的萨哈方程：
\begin{equation}\tag{9.8.10}\label{eq:9.8.10}
f _ { \mathrm { H } } = \frac { n _ { H } } { n _ { p } + n _ { \mathrm { H } } } = ( 1 - f _ { \mathrm { H } } ) ^{2} s ,
\end{equation}
其中参数 \(s\) 为
\begin{equation}\tag{9.8.11}\label{eq:9.8.11}
s = 4 \zeta ( 3 ) \sqrt { \frac { 2 } { \pi } } ( 1 - 2 q ) \eta \left( \frac { k T } { m _ { e } c ^{2} } \right) ^{3 / 2} e ^{\beta \mathrm { R y} } .
\end{equation}
\autoref{eq:9.8.10} 的解，即
\begin{equation}\tag{9.8.12}\label{eq:9.8.12}
f _ { \mathrm { H } } = \frac { 1 + 2 s - \sqrt { 1 + 4 s } } { 2 s } ,
\end{equation}
如图~9.8所示。在高于再结合温度的温度下，\(s\) 较小，因此 \(f_{\mathrm{H}}\) 较小，导致等离子体完全电离。而在低温下，\(s\) 较大，\(f_{\mathrm{H}}\) 接近于1，仅留下中性原子。重子与光子比值 \(\eta\) 和 \(Ry/m_{e}c^{2}\) 的数值较小，使得再结合在温度
\begin{equation}\tag{9.8.13}\label{eq:9.8.13}
k T _ { r } \approx \frac { \mathrm { R y } } { \ln \left( \frac { 1 } { \eta } \left( \frac { m _ { e } c ^{2} } { \mathrm { R y } } \right) ^{3 / 2} \right) } \, ,
\end{equation}
这使得最后一次散射推迟到 \(T \approx 3000\) K；详见图~9.8。
9.2. 计算宇宙最初1秒内，光子、电子、正电子和中微子的平均粒子能量以及平均粒子熵。
9.3. 尽管在早期宇宙中，光子与带电电子和正电子之间的电磁相互作用使它们彼此保持平衡，但请说明电子与正电子之间的直接电磁库仑相互作用能相较于这些粒子的相对论动能而言十分微小。请展示库仑能与动能之比大致等于精细结构常数：
\begin{equation}\tag{9.8.14}\label{eq:9.8.14}
\frac { u _ { \mathrm { c o u l o m b } } } { u _ { e } } \approx \alpha = \frac { e ^{2} } { 4 \pi \epsilon _ { 0 } \hbar c } = \frac { 1 } { 13 7 . 03 6 } .
\end{equation}
9.4. 请证明，在电子-正电子湮灭的早期阶段，电子数密度与光子数密度之比随温度变化的规律为：
\begin{equation}\tag{9.8.15}\label{eq:9.8.15}
\frac { n _ { - } } { n _ { \gamma } } \approx \frac { n _ { + } } { n _ { \gamma } } \sim \left( \frac { k T } { m _ { e } c ^{2} } \right) ^{3 / 2} \exp \Bigl ( - \beta m _ { e } c ^{2} \Bigr ) .
\end{equation}
9.5. 请证明，在几乎所有的正电子都已湮灭、电子数密度已接近与质子数密度趋于平稳之后，正电子数密度与光子数密度之比仍随温度变化的规律为：
\begin{equation}\tag{9.8.16}\label{eq:9.8.16}
\frac { n _ { + } } { n _ { \gamma } } \sim \left( \frac { k T } { m _ { e } c ^{2} } \right) ^{3 / 2} \exp  \left( - 2 \beta m _ { e } c ^{2} \right) .
\end{equation}
9.6. 在正电子湮灭之后，宇宙的能量密度主要由光子和中微子所主导。证明在那个时期，宇宙的能量密度为：\(u_{\mathrm{total}}=(1+(4/11)^{4/3})u_{\gamma}\)。接下来，利用哈勃膨胀关系\autoref{eq:9.1.1}、温度标度关系\autoref{eq:9.3.17}，以及电子-正电子湮灭后的能量密度，证明光子温度随时间的变化规律为：\(T(t)\approx10^{10}\,\mathrm{K}\sqrt{\frac{1.788\,\mathrm{s}}{t}}\)。这一关系从\(t\approx100\,\mathrm{s}\)一直持续到\(t\approx200,000\)年，此时由重子物质和冷暗物质所引起的能量密度开始占据主导地位。
3.7. 如果当前宇宙背景辐射的温度为27K，而非2.7K，宇宙的原始氦含量会受到怎样的影响？如果温度为0.27K又会怎样？
9.8. 在布鲁克海文国家实验室的相对论重离子对撞机（RHIC）中，金对金核碰撞会产生一种夸克-胶子等离子体，其能量密度约为\(4\,\mathrm{GeV}/\mathrm{fm}^3\)；参见Adare等人（2010）。将核物质视为由非相互作用的相对论性夸克-胶子气体构成。纳入低质量的上夸克和下夸克及其反粒子（均为自旋-\(\frac{1}{2}\)），以及自旋为1的无质量胶子。与光子一样，胶子是玻色子，每种胶子具有两种自旋态，并且自身就是其反粒子。胶子共有八种类型，能够改变夸克的三种颜色状态。由于等离子体的尺寸极小，只需考虑强相互作用的粒子。那么，夸克-胶子等离子体的温度是多少？
9.9. 计算宇宙早期，当温度高于\(kT = 300\,\mathrm{MeV}\)时，能量密度与温度之间的关系。在这些温度下，夸克和胶子已从单个原子核中被释放出来。将夸克-胶子等离子体视为非相互作用的相对论性气体。在那时，
在这些温度下，彼此处于平衡状态的物种包括：光子、三种中微子、电子和正电子、μ子和反μ子、上夸克和下夸克及其反粒子（均为自旋-\(\frac{1}{2}\)），以及自旋为1的无质量胶子。与光子一样，胶子是玻色子，每种胶子具有两种自旋态，并且自身就是其反粒子。胶子共有八种类型，能够改变夸克的三种颜色状态。奇异夸克、粲夸克、顶夸克和底夸克以及τ轻子的质量均超过300MeV，因此在这一温度下它们的贡献并不显著。利用你的计算结果和\autoref{eq:9.1.1}，确定该时期宇宙年龄与\(kT \approx 300\,\mathrm{MeV}\)时的年龄之间，温度随时间的变化规律。


